\section{Retransmit or Backfill?}
\label{sec:handling}

The InWaterSense and Weija records were replayed as the link state for NB-IoT (search times \nbSearchTimes~s), LoRaWAN SF7, and LoRaWAN SF12, with and without independent radio losses: \nConfigs{} configurations in all. Table~\ref{tab:strategies} shows InWaterSense under the full-slot model; ranges are over all configurations, and further results are given in Supplementary Section~\ref{sec:s-handling}.

\begin{table*}[t]
\centering
\caption{Loss-handling strategies on the InWaterSense record (full-slot model, no added radio loss): transmission energy per day, $E_\mathrm{tx}$ in (\ref{eq:energy}), per radio (NB-IoT: $T_\mathrm{search}=10$~s), and the share of readings delivered and delivered within 1~h (LoRaWAN SF7; NB-IoT for the NB-IoT-only strategy). ARQ: LoRaWAN 1.0.2 behaviour. Back-off (\ref{eq:backoff}) capped at 6~h. --: not applicable.}
\label{tab:strategies}
\small
\begin{tabular}{@{}lrrrrr@{}}
\toprule
 & \multicolumn{3}{c}{Energy (J/day)} & \multicolumn{2}{c}{Readings} \\
\cmidrule(lr){2-4}\cmidrule(l){5-6}
Strategy & NB-IoT & SF7 & SF12 & Delivered & Within 1 h \\
\midrule
\tabin{tables/tab_strategies.tex}
\bottomrule
\end{tabular}
\end{table*}

\subsection{ARQ pays for losses it can hardly recover}
\label{sec:hand-arq}
Under the full-slot model, ARQ delivered exactly what unconfirmed reporting delivered, at \arqEightMin--\arqEightMax{} its energy on LoRaWAN with eight transmissions. This is the pessimistic end of a bracket. If the link may return during a gap's last slot, retries recovered \arqGainUniformMin--\arqGainUniformMax{} percentage points of InWaterSense's readings on LoRaWAN; if, in addition, gaps of up to an hour are radio-loss bursts, they recovered \arqGainBurstMin--\arqGainBurstMax{} points. Against the \arqLostPp{} points that unconfirmed reporting lost, ARQ thus recovered under a fifth in every model, at \arqRatioMin--\arqRatioMax{} the energy of unconfirmed reporting; on Weija, it recovered at most \arqGainWeijaMax{} points. NB-IoT ARQ recovered nothing because the modem reports that there is no network. LoRaWAN 1.0.4 (eight transmissions at a fixed rate) changed the price only slightly: \arqVOneZeroFour{} instead of \arqVOneZeroTwo{} at SF7. Retransmission suits the radio regime, not the long gaps that hold most of the missing data.

\subsection{Buffering buys completeness at a modest energy cost}
\label{sec:hand-buffer}
On NB-IoT, searching for a cell in every slot of a gap cost \backoffMin--\backoffMax{} the energy of backing off; thus, the fair NB-IoT baseline backs off. Against the fair baseline, confirmed reporting with a buffer and back-off delivered every generated reading at an energy change of \bufNbMin\% to \bufNbMax\% on NB-IoT and \bufLoMin\% to \bufLoMax\% on LoRaWAN. It was cheaper at SF7, where a confirmed success costs less than an unconfirmed report and the node stops transmitting into gaps, and more expensive at SF12. This completeness rests on three assumptions, each of which was tested:
\begin{enumerate}
\item The reading must exist: With up to half of the long-gap slots reassigned to a node that was off, buffering still delivered every generated reading, but only \offExpectedInws{} of the expected readings on InWaterSense.
\item The buffer must hold the backlog: The longest backlog was \maxBufferKB~KB; 2~KB of RAM lost \ramLossTwo{} of the readings and 64~KB none, and writing every reading to flash~\cite{winbond2018w25q128jv} cost at most \flashJDayMax~J/day.
\item The drain must leave room for live traffic: Limited to what the duty cycle leaves after one live frame per slot, the drain still delivered everything, but at SF12, a backlog persisted for \drainTwelveMaxD~days in InWaterSense's worst period; therefore, live readings should go ahead of the backlog~\cite{ingelrest2010sensorscope}.
\end{enumerate}

\subsection{Batching is cheaper and costs little timeliness}
\label{sec:hand-batch}
Hourly batching with buffer and back-off delivered every reading, while the fair per-reading baseline used \batchSaveMin--\batchSaveMax{} as much energy. Its share of readings within 1~h differed from that of unconfirmed reporting by only \batchWithinMin{} to \batchWithinMax{} points: under long gaps, the gaps, not the batch interval, decide how many readings arrive within the hour. The back-off cap trades energy against how long the server waits once the network is back. Probing every slot cost \kEverySlotSeven{} the energy of $K=6$~h at SF7, while raising $K$ to 24~h saved only \kSixToDaySave{} more and tripled the 90th-percentile wait (\kPninetySix{} to \kPninetyDay~h); $K=6$~h sits at this knee. The price of batching is the timeliness of alarms and of the real-time picture of fast parameters, which Section~\ref{sec:design} addresses.

\subsection{Many nodes and model uncertainty}
\label{sec:hand-network}
A system outage ends for every node behind the gateway simultaneously. A discrete-event model of one EU868 gateway and up to 500 nodes (collisions with capture, eight demodulators, a half-duplex gateway, and gateway duty cycles; Supplementary Table~\ref{tab:s-netsim}) showed two effects. Without randomisation, the deterministic back-off kept the nodes synchronised: from \netZeroJitterN{} nodes at SF7, almost none recovered within 72~h. With probes delayed randomly by up to an hour, every SF7 node recovered, the last within \netSevenAllMax~h, at up to \netSevenEnergyMax{} the energy of a node alone. At SF12, the gateway's acknowledgements ran out: its first-window sub-band reached its 1\% limit, up to \netTwelveNoAck{} of the received frames went unacknowledged, and after a one-day outage, only \netTwelveCleared{} of 100 nodes recovered within 72~h. The per-node results hold for cells with tens of nodes.

With every measured energy component multiplied by an independent factor between 1/2 and 2 (\mcDraws{} draws), ARQ cost more than unconfirmed reporting on LoRaWAN in \mcArqDearer{} of the draws. Hourly batching was cheaper than the fair baseline in at least \mcBatchCheapOther{} of them at SF7 and on NB-IoT, but only in \mcBatchCheapTwelveMin--\mcBatchCheapTwelveMax{} at SF12. The buffering energy relative to the fair baseline ranged from \mcBufLo{} to \mcBufHi{}, and its sign is set by the ratio of confirmed to unconfirmed energy, which can only be settled by measuring the target hardware.

In response to RQ2, ARQ recovers only the short radio losses under any timing model, and on LoRaWAN, it multiplies energy. Buffering with back-off recovers every generated reading, given tens of kilobytes of storage, at an energy close to that of recovering nothing in the model, and hourly batching is the energy-efficient way to do it at SF7 and on NB-IoT. Many nodes recovering simultaneously need randomised probes, and large SF12 cells are limited by the gateway's acknowledgements.
